\section{Introducción}

\subsection{Propósito}
Este documento define los requerimientos técnicos del rover Olympus para una misión lunar de exploración, así como los criterios de validación del sistema integrado. Este documento está dirigido al equipo de ingeniería, asesores de misión y partes interesadas en el desarrollo del rover.

\subsection{Alcance}
El sistema Olympus es un rover autónomo diseñado para:
\begin{itemize}
    \item Exploración de superficie lunar.
    \item Navegación autónoma en terrenos no estructurados.
    \item Recolección y transmisión de datos científicos.
    \item Operación en entornos análogos terrestres para validación.
\end{itemize}

\subsection{Definiciones y Acrónimos}
\begin{table}[H]
    \centering
    \begin{tabularx}{\textwidth}{l X}
        \toprule
        \textbf{Acrónimo} & \textbf{Definición} \\
        \midrule
        GNC & Guidance, Navigation \& Control (Guía, Navegación y Control). \\
        RTOS & Real-Time Operating System (Sistema Operativo en Tiempo Real). \\
        TRL & Technology Readiness Level (Nivel de Madurez Tecnológica). \\
        ISRU & In-Situ Resource Utilization (Utilización de Recursos in situ). \\
        ROS 2 & Robot Operating System 2 (Middleware de robótica). \\
        MTBF & Mean Time Between Failures (Tiempo medio entre fallos). \\
        \bottomrule
    \end{tabularx}
\end{table}

\subsection{Referencias}
\begin{itemize}
    \item ISO/IEC/IEEE 29148:2018 - Systems and software engineering — Life cycle processes — Requirements engineering.
    \item NASA Systems Engineering Handbook.
    \item ECSS Standards (European Cooperation for Space Standardization).
\end{itemize}
